\documentclass{article}
\usepackage{amsmath,amssymb}
\usepackage{xurl}
\usepackage[hidelinks]{hyperref}
\setlength{\emergencystretch}{2em}
% Shared commands for notes in docs/paper-gaps/.

\DeclareUnicodeCharacter{2080}{\ifmmode{}_{0}\else\textsubscript{0}\fi}
\DeclareUnicodeCharacter{2081}{\ifmmode{}_{1}\else\textsubscript{1}\fi}
\DeclareUnicodeCharacter{2082}{\ifmmode{}_{2}\else\textsubscript{2}\fi}
\DeclareUnicodeCharacter{2099}{\ifmmode{}_{n}\else\textsubscript{n}\fi}

\newcommand{\Lean}{\textsc{Lean}}
\ExplSyntaxOn
\NewDocumentCommand{\leanid}{m}{
  \ifmmode
    \textsf{\detokenize{#1}}
  \else
    \group_begin:
    \urlstyle{sf}
    \tl_set:Nn \l_tmpa_tl {#1}
    \tl_replace_all:Nnn \l_tmpa_tl {\_} {_}
    \exp_args:NV \path \l_tmpa_tl
    \group_end:
  \fi
}
\ExplSyntaxOff
\providecommand{\tr}{\operatorname{tr}}
\newcommand{\tnleanIssueBase}{https://github.com/LionSR/TNLean/issues/}
\newcommand{\tnleanPrBase}{https://github.com/LionSR/TNLean/pull/}
\newcommand{\tnleanDocBase}{https://sirui-lu.com/TNLean/docs/}
\newcommand{\ghissue}[1]{\href{\tnleanIssueBase#1}{GitHub issue \##1}}
\newcommand{\ghpr}[1]{\href{\tnleanPrBase#1}{PR \##1}}
\newcommand{\doclink}[1]{\href{\tnleanDocBase#1.html}{\path{#1}}}

% Dirac notation and the identity operator, as in blueprint/src/macros/common.tex.
% \providecommand keeps note-local definitions authoritative.
\usepackage{dsfont}
\providecommand{\ket}[1]{|#1\rangle}
\providecommand{\bra}[1]{\langle #1|}
\providecommand{\braket}[2]{\langle #1 | #2 \rangle}
\providecommand{\ketbra}[2]{|#1\rangle\!\langle #2|}
\providecommand{\Id}{\mathds{1}}
\providecommand{\one}{\mathds{1}}
\providecommand{\C}{\mathbb{C}}
\providecommand{\MN}[1]{M_{#1}(\C)}
\providecommand{\MD}{M_D(\C)}

% Verdict marker: \gapnote{<kind>}{<status>}. Kind is the policy
% classification (clarification, local-correction, scope-restriction,
% unfaithful, false-source, open-gap); status is open, wip (elimination
% underway), resolved, or historical. Machine-read by the site index;
% prints a verdict line.
\newcommand{\gapnote}[2]{\par\noindent\textbf{Verdict:} #1 (#2).\par}

\title{The Block-Injective Interpolation: Action, Fusion, and the Ground Space}
\author{}
\date{2026-11-16}
\begin{document}
\maketitle
\gapnote{scope-restriction}{open}

\section{The source claim}
Garre-Rubio, Lootens and Moln\'ar extend the interpolation between two
symmetric states with equal $L$-symbols to states with a degenerate ground
space \cite[Section~5]{GarreRubio2023MPOSym}:
\begin{quote}
    ``We now deal with the case of systems with ground state degeneracy. We
    want to connect two block-injective MPSs (with the same number of
    blocks) that share the $L$-symbols under the action of two MPO algebras
    with the same fusion category $\mathcal{C}$ and same $F$-symbols.
    [\ldots] The embedding MPS tensor has the block-diagonal form
    $A(\gamma)=\bigoplus_x A_x(\gamma)$.''
\end{quote}
The section closes with
\begin{quote}
    ``The previous MPO tensors define an MPO representation of
    $\mathcal{C}$ if $L^0=L^1$ and $F^0=F^1$, this can be seen as in the
    injective case. Also, it corresponds to an MPO symmetry of
    $|\psi_{A(\gamma)}\rangle$ since it satisfies [\ldots] The injectivity
    of every block of $A(\gamma)$ and the injectivity of $T_a$ is justified
    in the same way as in the previous case. Moreover, $H(\gamma)$ is
    well-behaved in the whole path its ground space is spanned by
    $\{ |\psi_{A_x(\gamma)}\rangle\}$.''
\end{quote}
The injective case is introduced with the remark that ``the general case
with ground state degeneracy is a straightforward
generalization''.\footnote{Source traceability:
    \path{Papers/2203.12563/REsubmission.tex}, line 1582 (straightforward
    generalization), lines 1695--1704 (block correspondence and
    $A_x(\gamma)$), lines 1706--1726 (the common operators and their
    fusion), lines 1726--1777 (the action relation, injectivity, and the
    ground space).}

\section{Does the block argument use commutation?}
None of the assertions of lines 1695--1777 invokes commutation of the path
Hamiltonians $H(\gamma)$ with the periodic operators $O_a$: the passage
asserts the block-diagonal form of $A(\gamma)$, the definition of $T_a$,
its fusion under $L^0=L^1$ and $F^0=F^1$, the local action relation, the
injectivity of the blocks and of $T_a$, and the ground space of
$H(\gamma)$. The action and fusion statements formalized below therefore
need no adjoint closure and carry no restriction from
\texttt{glm23\_mixed\_path\_parent\_commutation.tex}.

The purpose of the section, to \emph{connect} the two states in the sense
of the definition of a phase, does require that $H(\gamma)$ commute with an
injective MPO representation. The block passage inherits this from the
injective case through line 1582. With one block the construction below is
exactly the injective interpolation, so the counterexample of
\texttt{glm23\_mixed\_path\_parent\_commutation.tex} (an involutive
non-unitary on-site symmetry whose canonical parent does not commute with
it) is also a block-injective counterexample. Any commutation statement for
the block path must carry the adjoint-closure restriction recorded there.
No such statement is made here.

\section{The formal boundary}
The blocks of both endpoints are indexed by one finite set, which records
the correspondence $x:x_0\to x_1$, and the action multiplicities are
common to both endpoints, which records the correspondence of the
degeneracies. The endpoint blocks are assumed to occupy separated physical
sectors, so that the local space of each endpoint is
$\bigoplus_x\mathbb C^{D^p_x}\otimes\mathbb C^{D^p_x}$. For several
blocks this is a restriction beyond the source: lines 1584--1586 restrict
the local space only to the span of the whole tensor, and the separation of
that span into block sectors is not derived (see the last paragraph of this
section).

The common operators act on every block by an exact biorthogonal
decomposition into the blocks $A_y(\gamma)$, for every real $\gamma$,
without injectivity, fusion, or equal $L$-symbols. The periodic operators
map each block state to the multiplicity-weighted sum of block states.
Each block is continuous and injective on $(0,1)$. When all endpoint blocks
are injective with positive dimension and the chosen raw $L$ matrices agree
for every pair of labels and every pair of corresponding blocks, the common
operators fuse exactly with the original multiplicities and satisfy the
fusion rules at every length. The common $F$-symbols are not needed for
exact fusion.\footnote{\leanid{MPSTensor.MPOSymmetry.isBiorthogonalDecomposition_blockMixedInterpolation},
    \leanid{MPSTensor.MPOSymmetry.mpo_mulVec_blockMixedInterpolation},
    \leanid{MPSTensor.MPOSymmetry.isInjective_blockMixedInterpolation},
    in \doclink{TNLean/MPS/Symmetry/MPOSymmetry/MixedPath/MPO/BlockMixedMPO}, and
    \leanid{MPSTensor.MPOSymmetry.isBiorthogonalDecomposition_blockMixedMPO},
    \leanid{MPSTensor.MPOSymmetry.isMPOFusionAlgebra_blockMixedMPO},
    in \doclink{TNLean/MPS/Symmetry/MPOSymmetry/MixedPath/MPO/BlockMixedMPOFusion}.}

The ground-space claim is proved for the two-site canonical parent
Hamiltonian $H(\gamma)$ of $A(\gamma)$, the Hamiltonian of the injective
case (line 1691). Since the blocks occupy separated physical sectors and
each block is injective on $(0,1)$, the one-site letters of the blocks
already span $\bigoplus_x M_{D^0_x+D^1_x}$, so the block-injective parent
theory applies with the source's two-site interaction. For every
$\gamma\in(0,1)$ and every ring of length $N\ge2$ the periodic ground space
of $H(\gamma)$ is spanned by the block states
$\ket{\psi_{A_x(\gamma)}}$, and these are linearly independent on the
whole closed interval, because they are nonzero and supported on disjoint
sectors. The statement that $H(\gamma)$ is well behaved is proved in the
following sense: the ground projection is continuous on $(0,1)$, every
compact subset of $(0,1)$ has one positive gap bound for all $N\ge2$, and
as $\gamma$ tends to any point of $[0,1]$, in particular to an endpoint,
the ground projections converge to the projection onto the span of the
block states at that point.\footnote{%
    \leanid{MPSTensor.MPOSymmetry.ker_hamiltonian_parentRingInteraction_blockMixedTensor},
    \leanid{MPSTensor.MPOSymmetry.linearIndependent_mpv_blockMixedInterpolation},
    \leanid{MPSTensor.MPOSymmetry.continuous_blockMixed_groundSpace_starProjection},
    \leanid{MPSTensor.MPOSymmetry.exists_uniform_blockMixed_hamiltonian_parentRingInteraction_gap},
    \leanid{MPSTensor.MPOSymmetry.tendsto_blockMixed_groundSpace_starProjection},
    in \doclink{TNLean/MPS/Symmetry/MPOSymmetry/MixedPath/MPO/BlockMixedGroundSpace}.}
All of these statements carry the separated-sector restriction.

The one-site injectivity of the operators $T_a$ on the enlarged space is
false already for one block, by the counterexample of
\texttt{glm23\_mixed\_operator\_one\_site\_injectivity.tex}. Open are the
identification of the limiting projections with the kernels of the
endpoint Hamiltonians, and a gap there, which the injective case proves
separately; and the restriction of
an arbitrary physical alphabet to the span of a block-injective tensor,
derived for one block in the arbitrary-alphabet construction but not for
several blocks. The source gives no argument for the ground-space claims
beyond the injective case.

\section{Elimination plan}
\begin{enumerate}
    \item The interior ground space, continuity, and compact gaps are
          proved. Identify the limiting projections at $\gamma=0,1$ with the
          kernels of the endpoint Hamiltonians, and prove an endpoint gap,
          as in the injective case.
    \item State any commutation of $H(\gamma)$ with the common operators
          under the adjoint-closure restriction of
          \texttt{glm23\_mixed\_path\_parent\_commutation.tex}, or replace
          the canonical parent by the symmetrized parent.
    \item Derive the restriction to the span of the tensor for several
          blocks with arbitrary alphabets, from a simultaneous span of the
          blocks, as in the one-block polar construction.
\end{enumerate}

\bibliographystyle{alpha}
\bibliography{references}
\end{document}
